\documentclass[11pt]{article}

\usepackage[a4paper,margin=2.4cm]{geometry}
\usepackage{amsmath,amssymb,amsthm}
\usepackage{mathtools}
\usepackage{enumitem}
\usepackage[colorlinks=true,linkcolor=blue,urlcolor=blue]{hyperref}

\newcommand{\R}{\mathbb{R}}
\newcommand{\C}{\mathbb{C}}
\newcommand{\E}{\mathbb{E}}
\newcommand{\Real}{\operatorname{Re}}
\newcommand{\ii}{\mathrm{i}}
\newcommand{\dd}{\,\mathrm{d}}
\newcommand{\ch}{\varphi}
\DeclarePairedDelimiter{\abs}{\lvert}{\rvert}

\newtheorem{remark}{Remark}

\title{\textbf{Lecture 6 --- Fourier methods for option pricing}\\[2pt]
\large Study notes: the Fourier family, density recovery by COS, and the COS pricing formula}
\author{Computational Finance (WI4151), TU Delft --- F.\ Fang}
\date{}

\begin{document}
\maketitle
\vspace{-1.2em}

\begin{abstract}
\noindent These notes work through Lecture 6 in full. The destination is the \emph{COS pricing
formula} for a European option,
\[
P(x_0,T)\;\approx\; e^{-rT}{\sum_{k=0}^{N-1}}' \Real\!\left\{\ch_{X_T}\!\Big(\tfrac{k\pi}{b-a};x_0\Big)\,
e^{-\ii k\pi\frac{a}{b-a}}\right\}V_k ,
\]
and everything needed to reconstruct it on paper: the characteristic function as the Fourier dual of the
density, why a \emph{cosine} (not full complex) series, the closed form of the cosine coefficients $A_k$ in
terms of the characteristic function $\ch$, the closed form of the payoff coefficients $V_k$ for a call/put,
the cumulant-based truncation range $[a,b]$, and the error analysis. The closed forms of $V_k$ and of $[a,b]$
are not on the slides; they are taken from Fang \& Oosterlee (2008/2009) and derived here because they are
the parts an examiner will press on. The prime on the sum means the $k=0$ term is halved.
\end{abstract}

\tableofcontents

\section{Where COS sits, and what problem it solves}

The risk-neutral price of a European-style claim with maturity $T$, payoff $V(x_T,T)$ and discount rate $r$ is
the discounted expectation
\begin{equation}\label{eq:pricing}
P(x_0,T)=e^{-rT}\,\E\!\left[V(x_T,T)\mid x_0\right]
= e^{-rT}\int_{\R} V(y,T)\,f(y\mid x_0)\dd y ,
\end{equation}
where $x_T$ is the state variable (typically $x_T=\ln S_T$ or the log-return $\ln(S_T/S_0)$) and
$f(\,\cdot\mid x_0)$ is its risk-neutral conditional density. The numerical problem is that $f$ is rarely
known in closed form. What \emph{is} known, for an enormous class of models (Black--Scholes, Heston,
Variance-Gamma, the affine jump-diffusions, the L\'evy models tabulated in \S\ref{sec:cf}), is the
\emph{characteristic function} $\ch_{X_T}$. COS is a way to compute the integral~\eqref{eq:pricing} using only
$\ch$, with no numerical integration and with (for smooth densities) exponential convergence in the number of
terms. It is the method this course is built around.

The logic of the lecture is three-fold and these notes follow it:
\begin{enumerate}[leftmargin=*]
  \item \textbf{The Fourier family} (\S\ref{sec:family}): the continuous Fourier transform, the characteristic
  function as a special case, the inversion formula, the DFT/FFT, and Fourier(-cosine) series. This fixes the
  vocabulary and the one identity COS exploits.
  \item \textbf{Density recovery by COS} (\S\ref{sec:cos-density}): approximate a density on a finite range
  $[a,b]$ by a cosine series whose coefficients are \emph{sampled directly from $\ch$}.
  \item \textbf{COS for option pricing} (\S\ref{sec:cos-pricing}): insert the cosine expansion into
  \eqref{eq:pricing}, swap sum and integral, and read off the pricing formula; then compute $V_k$ in closed
  form and fix the range $[a,b]$ from cumulants.
\end{enumerate}

\section{The Fourier family}\label{sec:family}

\subsection{Continuous Fourier transform (CFT)}
For integrable $f:\R\to\C$ one convention of the transform / inverse pair is
\begin{equation}\label{eq:cft}
\hat f(\omega)=\int_{-\infty}^{\infty} f(x)\,e^{\ii 2\pi\omega x}\dd x,
\qquad
f(x)=\frac{1}{2\pi}\int_{-\infty}^{\infty}\hat f(\omega)\,e^{-\ii 2\pi\omega x}\dd\omega .
\end{equation}
$\hat f(\omega)$ is the amplitude/phase of the frequency-$\omega$ component; the inversion rebuilds $f$ as a
superposition of complex exponentials of all frequencies. (Conventions differ in the placement of $2\pi$ and
the sign in the exponent; only internal consistency matters.)

\subsection{The characteristic function as Fourier dual of the density}\label{sec:cf}
The characteristic function of a random variable $X$ is
\begin{equation}\label{eq:chf}
\ch_X(t)=\E\!\left[e^{\ii t X}\right]=\int_{-\infty}^{\infty} e^{\ii t x} f_X(x)\dd x .
\end{equation}
Comparing \eqref{eq:chf} with \eqref{eq:cft}, $\ch_X$ \emph{is} (up to the $2\pi$-scaling convention) the
Fourier transform of the density $f_X$. The CDF $F_X$ determines $\ch_X$ and vice versa; when a density exists,
$f_X$ and $\ch_X$ are a Fourier pair. The inversion reads
\begin{equation}\label{eq:inversion}
f_X(x)=\frac{1}{2\pi}\int_{-\infty}^{\infty} e^{-\ii t x}\,\ch_X(t)\dd t .
\end{equation}
\emph{This is the integral COS evaluates cheaply.} Before COS, \eqref{eq:inversion} (or the equivalent pricing
integral) was attacked by direct quadrature or by FFT (Carr--Madan, \S\ref{sec:carrmadan}).

\begin{remark}[Why the characteristic function is the right object]
The whole appeal is in \S\ref{sec:cf}: a long list of pricing models give $\ch$ in closed (or semi-closed)
form even when $f$ is intractable --- Black--Scholes, Heston, Variance-Gamma, Merton/Kou jumps, the affine
jump-diffusions, CGMY, NIG. Some of these are tabulated in \S\ref{sec:cftable}. A method that needs only $\ch$
inherits this entire library for free.
\end{remark}

\subsection{Discrete Fourier transform (DFT) and the FFT}
On a sequence $\{x_n\}_{n=0}^{N-1}$ the DFT and its inverse are
\begin{equation}
X_k=\sum_{n=0}^{N-1} x_n\,e^{\ii\frac{2\pi}{N}kn},
\qquad
x_n=\frac{1}{N}\sum_{k=0}^{N-1} X_k\,e^{-\ii\frac{2\pi}{N}kn}.
\end{equation}
Computed naively this is $O(N^2)$. The \emph{Fast Fourier Transform} factorises the DFT matrix into
$O(\log N)$ sparse factors (the radix-2 split into even/odd indexed sub-DFTs of length $N/2$), giving
$O(N\log N)$. The FFT is what makes the Carr--Madan approach practical and is the benchmark COS is measured
against. Note for later: COS does \emph{not} use the FFT --- its cost is the bare $O(N)$ evaluation of the
single sum.

\subsection{Fourier series and the half-range cosine series}\label{sec:halfrange}
A $2\pi$-periodic function expands as
$f(x)=\tfrac{a_0}{2}+\sum_{n\ge 1}\big(a_n\cos nx+b_n\sin nx\big)$. Restrict to a function on $[0,L]$ and extend
it to $[-L,L]$:
\begin{itemize}[leftmargin=*]
  \item the \textbf{even} extension kills all sine terms and gives the \emph{half-range cosine series}
  \begin{equation}\label{eq:halfcos}
  f(x)=\frac{a_0}{2}+\sum_{n=1}^{\infty} a_n\cos\!\Big(\frac{n\pi}{L}x\Big),
  \qquad a_n=\frac{2}{L}\int_0^L f(x)\cos\!\Big(\frac{n\pi}{L}x\Big)\dd x ;
  \end{equation}
  \item the \textbf{odd} extension kills all cosine terms and gives the half-range sine series.
\end{itemize}
COS uses the cosine version on a finite truncation interval $[a,b]$ (shift $x\mapsto x-a$, period $2(b-a)$).
The reason for choosing cosine over the full complex series is taken up in \S\ref{sec:whycos}.

\section{The first Fourier pricing method: Carr--Madan (1998)}\label{sec:carrmadan}

This is the precursor COS improves on; understanding it sharpens the contrast. Write $k=\ln K$,
$x_T=\ln S_T$. The call is
\begin{equation}
C_T(k)=e^{-rT}\int_{k}^{\infty}\big(e^{x_T}-e^{k}\big) f(x_T\mid x_0)\dd x_T .
\end{equation}
As $k\to-\infty$, $C_T(k)\to S_0$, so $C_T$ is not square-integrable in $k$ and cannot be Fourier-transformed
directly. Carr--Madan \emph{damp} it,
\begin{equation}
c_T(k)=e^{\alpha k}C_T(k),\qquad \alpha>0,
\end{equation}
so $c_T\in L^1\cap L^2$. Its Fourier transform $\psi_T(v)=\int_{-\infty}^{\infty} e^{\ii v k} c_T(k)\dd k$ can be
written in closed form in terms of the characteristic function:
\begin{equation}\label{eq:carrmadan}
\psi_T(v)=\frac{e^{-rT}\,\ch_X\!\big(v-(\alpha+1)\ii\big)}{\alpha^2+\alpha-v^2+\ii(2\alpha+1)v},
\end{equation}
and the price is recovered by an \emph{inverse} transform, evaluated by FFT:
\begin{equation}
C_T(k)=\frac{e^{-\alpha k}}{2\pi}\int_{-\infty}^{\infty} e^{-\ii v k}\,\psi_T(v)\dd v .
\end{equation}
Two structural weaknesses, both of which COS removes: (i) it still performs a numerical \emph{integration}
(the FFT computes a trapezoidal-type quadrature of an oscillatory integrand), and (ii) it needs a free damping
parameter $\alpha$, whose choice affects accuracy. COS needs neither.

\section{Density recovery by the COS method}\label{sec:cos-density}

\subsection{The cosine expansion and its coefficients}
Let $f$ be (essentially) supported on a finite interval $[a,b]$. Its half-range cosine series
\eqref{eq:halfcos}, written with the prime convention, is
\begin{equation}\label{eq:cosexp}
f(x)={\sum_{k=0}^{\infty}}' A_k\,\cos\!\Big(k\pi\frac{x-a}{b-a}\Big),
\qquad
A_k=\frac{2}{b-a}\int_{a}^{b} f(x)\cos\!\Big(k\pi\frac{x-a}{b-a}\Big)\dd x ,
\end{equation}
where $\sum'$ halves the $k=0$ term (so $A_0/2$ is the mean of $f$ over $[a,b]$).

\subsection{The key identity: $A_k$ from the characteristic function}\label{sec:Akderiv}
This is the engine of the method. Start from the coefficient integral and write the cosine as the real part
of a complex exponential:
\begin{align}
A_k&=\frac{2}{b-a}\int_a^b f(x)\cos\!\Big(k\pi\frac{x-a}{b-a}\Big)\dd x
   =\frac{2}{b-a}\,\Real\!\left\{\int_a^b f(x)\,
       \exp\!\Big(\ii k\pi\frac{x-a}{b-a}\Big)\dd x\right\}\notag\\
   &=\frac{2}{b-a}\,\Real\!\left\{e^{-\ii k\pi\frac{a}{b-a}}
       \int_a^b f(x)\,\exp\!\Big(\ii \frac{k\pi}{b-a}\,x\Big)\dd x\right\}.\label{eq:Akmid}
\end{align}
Now \emph{extend the integration range to $\R$}. Because $f$ has negligible mass outside $[a,b]$,
\begin{equation}\label{eq:trunc}
\int_a^b f(x)\,e^{\ii \frac{k\pi}{b-a}x}\dd x
\;\approx\;
\int_{-\infty}^{\infty} f(x)\,e^{\ii \frac{k\pi}{b-a}x}\dd x
\;=\;\ch_X\!\Big(\frac{k\pi}{b-a}\Big),
\end{equation}
the characteristic function~\eqref{eq:chf} evaluated at $t=k\pi/(b-a)$. Substituting \eqref{eq:trunc} into
\eqref{eq:Akmid} gives the coefficient \emph{without any integral}:
\begin{equation}\label{eq:Fk}
\boxed{\;A_k\approx F_k:=\frac{2}{b-a}\,\Real\!\left\{\ch_X\!\Big(\frac{k\pi}{b-a}\Big)
\exp\!\Big(-\ii k\pi\frac{a}{b-a}\Big)\right\}.\;}
\end{equation}
The single approximation made so far is the range truncation in \eqref{eq:trunc}; everything else is exact.
\emph{This is precisely why COS needs only $\ch$ and not $f$.}

\paragraph{Worked check: the standard normal.}
For $X\sim\mathcal N(0,1)$, $\ch_X(t)=e^{-t^2/2}$. With $a=-b$ symmetric,
\[
F_k=\frac{2}{b-a}\,\Real\!\left\{e^{-\frac12\left(\frac{k\pi}{b-a}\right)^2}
e^{-\ii k\pi\frac{a}{b-a}}\right\},
\]
matching the slide formula. A natural range is $a=\Phi^{-1}(10^{-8})$, $b=-a$ (clip the tails at probability
$10^{-8}$). With $K\!=\!50$ terms the reconstruction is graphically indistinguishable from the true density,
and at a fixed point the error falls \emph{exponentially} in $K$ until it saturates at the (fixed) range-truncation
floor --- the canonical COS convergence picture (\S\ref{sec:convergence}).

\subsection{Why a cosine series and not a full complex Fourier series}\label{sec:whycos}
Three reasons, each worth stating explicitly because it is a favourite question:
\begin{enumerate}[leftmargin=*]
  \item \textbf{Real coefficients, real arithmetic.} The even extension of $f$ on $[a,b]$ kills every sine term;
  the surviving coefficients $A_k$ are real (\eqref{eq:Fk} is a real part). A complex series would carry twice
  the coefficients and complex bookkeeping for no gain.
  \item \textbf{Direct sampling from $\ch$.} The cosine coefficient is exactly the real part of $\ch$ evaluated
  at the grid frequency $k\pi/(b-a)$ (\eqref{eq:Fk}). The cosine basis is the one whose coefficients are
  ``read off'' from the characteristic function in one line.
  \item \textbf{Payoff side closes in elementary functions.} The payoff coefficients $V_k$
  (\S\ref{sec:Vk}) are integrals of (polynomial/exponential)$\times\cos$, which integrate to elementary
  closed forms $\chi_k,\psi_k$. With a complex exponential basis the call payoff coefficients would not be as
  clean.
\end{enumerate}
Equivalently: the cosine series is the discrete cosine transform analogue, and the whole method is a clever way
to pair the cosine coefficients of $f$ (got from $\ch$) against the cosine coefficients of $V$ (got in closed
form) --- see \S\ref{sec:cos-pricing}.

\section{COS for European option pricing}\label{sec:cos-pricing}

\subsection{Deriving the pricing formula}
Insert the cosine expansion \eqref{eq:cosexp} of the (truncated) density into the price \eqref{eq:pricing} and
restrict the integral to $[a,b]$:
\begin{equation}
P(x_0,T)\approx e^{-rT}\int_a^b V(y,T)
{\sum_{k=0}^{\infty}}' A_k(x_0)\cos\!\Big(k\pi\frac{y-a}{b-a}\Big)\dd y .
\end{equation}
Interchange sum and integral (justified by uniform convergence on $[a,b]$ for smooth $f$) and define the
\textbf{payoff cosine coefficients}
\begin{equation}\label{eq:Vkdef}
V_k:=\frac{2}{b-a}\int_a^b V(y,T)\cos\!\Big(k\pi\frac{y-a}{b-a}\Big)\dd y .
\end{equation}
Then
\begin{equation}\label{eq:AkVk}
P(x_0,T)\approx e^{-rT}\,\frac{b-a}{2}{\sum_{k=0}^{\infty}}' A_k(x_0)\,V_k .
\end{equation}
Finally replace $A_k$ by its characteristic-function form \eqref{eq:Fk} --- note
$\tfrac{b-a}{2}A_k=\Real\{\ch_X(\tfrac{k\pi}{b-a})e^{-\ii k\pi a/(b-a)}\}$ --- and truncate the series at $N$
terms:
\begin{equation}\label{eq:cosprice}
\boxed{\;P(x_0,T)\approx e^{-rT}{\sum_{k=0}^{N-1}}'
\Real\!\left\{\ch_{X_T}\!\Big(\frac{k\pi}{b-a};x_0\Big)\,
\exp\!\Big(-\ii k\pi\frac{a}{b-a}\Big)\right\}V_k .\;}
\end{equation}
This is the European COS formula. The characteristic function enters in exactly one place: as the cosine
coefficients of the density, $\Real\{\ch\,e^{-\ii k\pi a/(b-a)}\}$. The dependence on the spot $x_0$ lives
inside $\ch_{X_T}(\,\cdot\,;x_0)$; for L\'evy/BS models $\ch_{X_T}(u;x_0)=\ch_{\text{lvl}}(u)\,e^{\ii u x_0}$,
so a single set of $\{V_k\}$ prices the whole strike/spot grid --- only the cheap $e^{\ii u x_0}$ factor
changes. Cost is $O(N)$ per price; no integration, no FFT.

\subsection{Closed form of the payoff coefficients $V_k$}\label{sec:Vk}
Work in the log variable $y=\ln(S_T/K)$, so $S_T=K e^{y}$ and the strike sits at $y=0$. Introduce the two
building-block integrals on a sub-interval $[c,d]\subseteq[a,b]$, with $\omega_k:=\dfrac{k\pi}{b-a}$:
\begin{align}
\chi_k(c,d)&:=\int_c^d e^{y}\cos\!\big(\omega_k(y-a)\big)\dd y,\label{eq:chi}\\
\psi_k(c,d)&:=\int_c^d \cos\!\big(\omega_k(y-a)\big)\dd y.\label{eq:psi}
\end{align}

\paragraph{Derivation of $\chi_k$.}
Write the cosine as a real part and integrate the complex exponential exactly:
\begin{align}
\chi_k(c,d)&=\Real\!\left\{\int_c^d e^{(1+\ii\omega_k)y}e^{-\ii\omega_k a}\dd y\right\}
=\Real\!\left\{\frac{e^{-\ii\omega_k a}}{1+\ii\omega_k}\,
\Big[e^{(1+\ii\omega_k)y}\Big]_c^d\right\}.
\end{align}
Multiply numerator and denominator by $(1-\ii\omega_k)$ (so the denominator becomes $1+\omega_k^2$) and take
the real part of $e^{y}e^{\ii\omega_k(y-a)}(1-\ii\omega_k)=e^y\big[\cos\omega_k(y-a)+\omega_k\sin\omega_k(y-a)\big]
+\ii(\cdots)$. This gives the standard closed form
\begin{equation}\label{eq:chiclosed}
\chi_k(c,d)=\frac{1}{1+\omega_k^2}\Big[
\cos\!\big(\omega_k(d-a)\big)e^{d}-\cos\!\big(\omega_k(c-a)\big)e^{c}
+\omega_k\sin\!\big(\omega_k(d-a)\big)e^{d}-\omega_k\sin\!\big(\omega_k(c-a)\big)e^{c}\Big].
\end{equation}

\paragraph{Derivation of $\psi_k$.}
Elementary:
\begin{equation}\label{eq:psiclosed}
\psi_k(c,d)=
\begin{cases}
\dfrac{1}{\omega_k}\Big[\sin\!\big(\omega_k(d-a)\big)-\sin\!\big(\omega_k(c-a)\big)\Big]
=\dfrac{b-a}{k\pi}\Big[\sin\!\big(\omega_k(d-a)\big)-\sin\!\big(\omega_k(c-a)\big)\Big], & k\neq 0,\\[2ex]
d-c, & k=0.
\end{cases}
\end{equation}

\paragraph{Call and put.}
For the call, $V=K(e^{y}-1)^{+}$ has support $y\in[0,\infty)$, so on $[a,b]$ the relevant interval is $[0,b]$:
\begin{equation}\label{eq:Vcall}
V_k^{\text{call}}=\frac{2}{b-a}\,K\int_0^b\big(e^{y}-1\big)\cos\!\big(\omega_k(y-a)\big)\dd y
=\frac{2}{b-a}\,K\big[\chi_k(0,b)-\psi_k(0,b)\big].
\end{equation}
For the put, $V=K(1-e^{y})^{+}$ has support $y\in(-\infty,0]$, so the interval is $[a,0]$:
\begin{equation}\label{eq:Vput}
V_k^{\text{put}}=\frac{2}{b-a}\,K\int_a^0\big(1-e^{y}\big)\cos\!\big(\omega_k(y-a)\big)\dd y
=\frac{2}{b-a}\,K\big[\psi_k(a,0)-\chi_k(a,0)\big].
\end{equation}
These $V_k$ are computed once. Put--call parity gives a fast cross-check and, in practice, the put is
preferred when its density mass sits in a numerically gentler region; either way \eqref{eq:cosprice} returns
the price.

\subsection{Greeks at no extra cost}
Because $S_0$ (equivalently $x_0$) enters \eqref{eq:cosprice} only analytically through $\ch_{X_T}(\cdot;x_0)$,
differentiation passes through the finite sum. With $\ch_{X_T}(u;x_0)=\ch(u)e^{\ii u x}$, $x=\ln(S_0/K)$,
\begin{align}
\Delta=\frac{\partial P}{\partial S_0}&\approx e^{-rT}{\sum_{k=0}^{N-1}}'
\Real\!\left\{\ch\!\Big(\tfrac{k\pi}{b-a}\Big)e^{\ii k\pi\frac{x-a}{b-a}}\cdot
\frac{\ii k\pi}{b-a}\right\}\frac{V_k}{S_0},\\
\text{Vega}=\frac{\partial P}{\partial u_0}&\approx e^{-rT}{\sum_{k=0}^{N-1}}'
\Real\!\left\{\frac{\partial\ch\!\big(\tfrac{k\pi}{b-a};u_0\big)}{\partial u_0}\,
e^{\ii k\pi\frac{x-a}{b-a}}\right\}V_k ,
\end{align}
where $u_0$ is the relevant model parameter (e.g.\ the initial variance in Heston). No re-integration is
needed --- a structural advantage over quadrature- and FFT-based pricers.

\section{The truncation range $[a,b]$ from cumulants}\label{sec:range}

The only approximation in the density step was discarding the tails outside $[a,b]$ in \eqref{eq:trunc}. The
range must therefore be wide enough to contain essentially all the probability mass of $X_T$, but no wider
(too wide wastes terms: a fixed accuracy needs $N\propto(b-a)$). Fang--Oosterlee fix it from the
\emph{cumulants} $c_n$ of $X_T$:
\begin{equation}\label{eq:abrange}
\boxed{\;[a,b]=\Big[c_1-L\sqrt{c_2+\sqrt{c_4}},\;\; c_1+L\sqrt{c_2+\sqrt{c_4}}\Big],
\qquad L\in[8,12]\ (\text{typ. }10).\;}
\end{equation}
Reading the terms:
\begin{itemize}[leftmargin=*]
  \item $c_1$ is the \textbf{mean} --- it centres the interval where the mass is;
  \item $c_2$ is the \textbf{variance} --- $\sqrt{c_2}$ is the natural width scale, so $L\sqrt{c_2}$ is an
  ``$L$-sigma'' box;
  \item $\sqrt{c_4}$ (from the fourth cumulant, controlling tail/kurtosis) \textbf{fattens} the box for
  heavy-tailed models (jumps, VG); for Gaussian models $c_4=0$ and it drops out;
  \item $L$ is the safety multiplier; larger $L$ lowers the truncation floor but demands more terms.
\end{itemize}

\paragraph{Black--Scholes cumulants.}
For $X_T=\ln(S_T/K)$ under GBM, $\ln S_T$ is normal, so only $c_1,c_2$ are nonzero:
\begin{equation}
c_1=\ln\tfrac{S_0}{K}+\Big(r-\tfrac12\sigma^2\Big)T,\qquad
c_2=\sigma^2 T,\qquad c_4=0,
\end{equation}
giving the clean box $[a,b]=c_1\pm L\sigma\sqrt{T}$.

\subsection{The short-maturity failure (and why it bites)}\label{sec:bug}
\emph{This is exactly the kind of trap an examiner sets, and a bug I hit in Assignment 2.} As $T\to0$,
$c_2=\sigma^2 T\to 0$, so the half-width $L\sqrt{c_2+\sqrt{c_4}}\to 0$ and $[a,b]$ \emph{collapses} towards the
single point $c_1$. Consequences:
\begin{itemize}[leftmargin=*]
  \item the interval can become so narrow that it fails to contain the region where the payoff $\times$ density
  product has mass --- in particular, in a Gaussian model with $c_4=0$ there is no $\sqrt{c_4}$ term to keep a
  floor, so the box shrinks unchecked;
  \item when pricing across a \emph{grid} of spots/strikes, a range computed for one $(S_0,K)$ may be centred
  by $c_1$ away from another point on the grid, so the truncated density is evaluated where the cosine
  reconstruction is meaningless --- producing large, erratic pricing errors precisely at short maturities;
  \item because the range error is a \emph{floor} on accuracy (\S\ref{sec:convergence}), simply adding more
  cosine terms $N$ does \emph{not} fix it --- the error plateaus.
\end{itemize}
\textbf{Fixes:} (i) increase $L$ at short maturity; (ii) keep the $\sqrt{c_4}$ term and/or impose a minimum
half-width floor so $b-a$ cannot collapse; (iii) make sure the range is computed to cover the whole
spot/strike grid actually priced, not a single point. The diagnostic signature is: errors that \emph{grow} as
$T\to 0$ and that \emph{do not} improve when $N$ is increased --- the tell-tale of a range, not a series,
problem.

\section{Error analysis and convergence}\label{sec:convergence}

The COS price \eqref{eq:cosprice} carries three error sources, and the lecture's central claim is that for
smooth densities the total error converges \emph{exponentially}.
\begin{enumerate}[leftmargin=*]
  \item \textbf{Range-truncation error} (from \eqref{eq:trunc}): the discarded tail mass outside $[a,b]$. It is
  \emph{fixed} once $[a,b]$ is chosen --- it sets the accuracy \emph{floor} at which the series convergence
  saturates. For densities with exponentially decaying tails it is itself exponentially small in the half-width.
  \item \textbf{Series-truncation error} (cutting the sum at $N$): governed by the decay of the cosine
  coefficients $A_k$ as $k\to\infty$. By the standard Fourier result, the smoother $f$, the faster $A_k$ decay:
  \begin{itemize}
    \item $f\in C^\infty$ with all derivatives integrable / $f$ analytic in a strip $\Rightarrow$
    \textbf{exponential} decay of $A_k$ $\Rightarrow$ exponential convergence in $N$;
    \item $f$ with only $\beta$ integrable derivatives (a finite-smoothness density --- e.g.\ near a jump, or
    VG with small $\nu$) $\Rightarrow$ only \textbf{algebraic} decay $A_k=O(k^{-(\beta+1)})$ $\Rightarrow$
    algebraic convergence. \emph{Lack of smoothness is what kills the exponential rate.}
  \end{itemize}
  \item \textbf{Characteristic-function / model error}: any error in $\ch$ itself (e.g.\ a mis-evaluated
  Heston $\ch$ across the branch cut) feeds straight through.
\end{enumerate}
The empirical picture (slide 24): at a fixed evaluation point the error drops like $e^{-cN}$ and then
\emph{plateaus} at the range-truncation floor --- raise the floor (wider $[a,b]$, larger $L$) and the plateau
drops, but you then need more terms to reach it. Compared with Monte Carlo's $O(N^{-1/2})$ and Carr--Madan/FFT,
COS reaches machine precision with $N$ in the dozens to low hundreds (the Heston tables on slide 31:
$N\!\sim\!160$ for $\sim\!10^{-10}$ error, versus thousands of FFT nodes for far worse error and longer CPU
time).

\section{COS versus Carr--Madan / FFT}\label{sec:vs}

\begin{center}
\renewcommand{\arraystretch}{1.25}
\begin{tabular}{p{4.3cm} p{5.0cm} p{5.0cm}}
\hline
& \textbf{COS} & \textbf{Carr--Madan / FFT} \\
\hline
Core operation & one $O(N)$ cosine sum, no integration & FFT quadrature, $O(N\log N)$ \\
Free parameters & none (range from cumulants) & damping $\alpha$ to be tuned \\
Convergence (smooth $f$) & exponential in $N$ & algebraic (quadrature-limited) \\
Terms for $\sim$machine prec. & dozens--hundreds & thousands of nodes \\
Greeks & analytic, from the same sum & extra transforms / finite differences \\
Needs only $\ch$? & yes & yes \\
\hline
\end{tabular}
\end{center}

\noindent\textbf{When COS wins:} smooth (or piecewise-smooth) densities with known $\ch$ and exponentially
decaying tails --- i.e.\ most European pricing under BS, Heston, and the standard L\'evy/AJD models. \textbf{When
to be careful:} densities with genuine low regularity (slow algebraic decay of $A_k$) and very short maturities
or fat tails, where $[a,b]$ must be set carefully (\S\ref{sec:bug}) and $N$ raised.

\section{Characteristic functions to keep at hand}\label{sec:cftable}

From Fang--Oosterlee (2009); here $\xi$ is the transform argument and $t$ the horizon.
\begin{center}
\renewcommand{\arraystretch}{1.6}
\begin{tabular}{l l}
\hline
\textbf{Model} & \textbf{Characteristic function} \\
\hline
Black--Scholes & $\ch(\xi,t)=\exp\!\big(\ii\xi\mu t-\tfrac12\sigma^2\xi^2 t\big)$ \\
NIG & $\ch(\xi,t)=\exp\!\big[\delta t\big(\sqrt{\alpha^2-\beta^2}-\sqrt{\alpha^2-(\beta+\ii\xi)^2}\big)\big]$ \\
Kou & $\ch(\xi,t)=\exp\!\big[\lambda t\big(\tfrac{p\eta_1}{\eta_1-\ii\xi}+\tfrac{(1-p)\eta_2}{\eta_2+\ii\xi}-1\big)\big]$ \\
Merton & $\ch(\xi,t)=\exp\!\big[\lambda t\big(\exp(\ii\bar\mu\xi-\tfrac12\bar\sigma^2\xi^2)-1\big)\big]$ \\
Variance-Gamma & $\ch(\xi,t)=\big(1-\ii\xi\theta\nu+\tfrac12\sigma^2\nu\xi^2\big)^{-t/\nu}$ \\
CGMY & $\ch(\xi,t)=\exp\!\big(Ct\,\Gamma(-Y)\big[(M-\ii\xi)^Y-M^Y+(G+\ii\xi)^Y-G^Y\big]\big)$ \\
\hline
\end{tabular}
\end{center}
\noindent For BS the risk-neutral drift is $\mu=r-\tfrac12\sigma^2$ in the log-price (martingale condition).
Heston's $\ch$ (Lecture 7) plugs into the same formula \eqref{eq:cosprice}; the Feller condition
$2\kappa\bar v\ge\eta^2$ keeps the variance process from hitting zero and keeps $\ch$ well-behaved.

\section{One-paragraph reconstruction (the thing to be able to say out loud)}

The price is a discounted expectation $e^{-rT}\!\int V f$. I do not know $f$, but I know its Fourier transform,
the characteristic function $\ch$. So I truncate to a finite range $[a,b]$ chosen from the cumulants
($c_1\pm L\sqrt{c_2+\sqrt{c_4}}$), expand $f$ there in a half-range \emph{cosine} series whose coefficients are
$A_k=\frac{2}{b-a}\Real\{\ch(\frac{k\pi}{b-a})e^{-\ii k\pi a/(b-a)}\}$ --- read straight off $\ch$, no integral.
I expand the payoff in the same cosine basis, getting coefficients $V_k$ that are elementary closed forms
($\chi_k,\psi_k$). Plancherel-style, the integral of the product becomes $\frac{b-a}{2}\sum' A_k V_k$, i.e.
$P\approx e^{-rT}\sum_{k=0}^{N-1}{}'\Real\{\ch(\frac{k\pi}{b-a})e^{-\ii k\pi a/(b-a)}\}V_k$. For smooth $f$ the
coefficients decay exponentially, so a few dozen terms reach machine precision; the only floor is the range
truncation, which is why a too-narrow $[a,b]$ at short maturity ($c_2=\sigma^2T\to0$) cannot be cured by adding
terms.

\end{document}
